\documentclass[12pt, titlepage]{article}
\usepackage{amsmath}
\usepackage{amsfonts}
\usepackage{amssymb}
\usepackage{fancyhdr}
\usepackage{unicode-math}
\AtBeginDocument{ \let\mathbb\relax \DeclareMathAlphabet{\mathbb}{U}{msb}{m}{n} }
\usepackage[mocha,styleAll]{catppuccinpalette}
\usepackage{cancel}
\usepackage{xcolor}
\usepackage[bidi=basic]{babel}
\babelprovide[main,import]{hebrew}
\babelfont[hebrew]{rm}{Hadasim CLM}
\babelfont[hebrew]{sf}{Miriam Mono CLM}
\usepackage{titlesec}
\newcommand\Ccancel[2][black]{\renewcommand\CancelColor{\color{#1}}\cancel{#2}}
\newcommand{\pard}[2][f]{\frac{\partial #1}{\partial #2}}
\title{חשבון אינפיניטסימלי 3 - 01040295 \\ גליון 1}
\author{יונתן אבידור - 214269565}
\titleformat{\section}{\Large\sffamily\bfseries}{\thesection}{1em}{}
\begin{document}
\maketitle
\pagestyle{fancy}
\fancyhead{}
\fancyfoot[L]{\text{214269565}}
\part*{שאלה 1}
כידוע כל הנורמות שקולות זו לזו ב-$\mathbb{R}^n$, ולכן לכל שתי נורמות $\lVert x\rVert_\alpha, \lVert x \rVert_\beta$ קיים קבוע $C_{\alpha,\beta}$ כך ש-
\[
  \left(\forall x\in\mathbb{R}^n\right) \left(  \lVert x \rVert_\alpha \leq C_{\alpha,\beta} \lVert x\rVert_\beta\right)
\]
עבור $\alpha,\beta=1,2,\infty$ יש להוכיח כי הקבוע $C_{\alpha,\beta}$ נתון ע"י הטבלה הבאה:
  \begin{center}
    \begin{tabular}[c]{c||c|c|c|}
      $_\alpha \setminus ^\beta$ & $1$ & $2$ & $\infty$\\
      \hline
      \hline
      $1$ & $*$ & $\sqrt{n}$ & $n$ \\
      \hline
      $2$ & $1$ & $*$ & $\sqrt{n}$ \\
      \hline
      $\infty$ & $1$ & $1$ & $*$\\
      \hline
    \end{tabular}
  \end{center}

לכל אחת מששת המקרים יש להראות כי $C_{\alpha,\beta}$ הנ"ל נותן את החסם הטוב ביותר, ע"י מציאת דוגמה לווקטור $x$ (שונה מאפס) שעבורו מתקבל שוויון.
\part*{פתרון 1}
יהי $x = \begin{pmatrix}x_1&\ldots&x_n\end{pmatrix}^t\in \mathbb{R}^n \neq 0_n$, נחפש $C_{\alpha,\beta}$ כך ש
\[
  \lVert x\rVert_\alpha \leq C_{\alpha,\beta} \lVert x \rVert_\beta 
\]
\section*{$C_{1,2}$}
נגדיר וקטור $y=\begin{pmatrix}\left|x_1\right|&\ldots&\left|x_n\right|\end{pmatrix}^t$, נשים לב כי $\lVert x\rVert_2 = \lVert y \rVert_2$.\\
בנוסף, נגדיר $z=\begin{pmatrix}1&\ldots&1\end{pmatrix}^t$, ונשים לב כי $\lVert z\rVert_2 = \sqrt{\sum_{i=1}^n 1^2} = \sqrt{n}$\\
לפי אי-שוויון קושי שוורץ עם הנורמה $\lVert \cdot \rVert_2$ (אפשר להשתמש איתה באי-שוויון קושי-שוורץ, שכן היא מושרית על ידי מכפלה פנימית) נקבל
\[
  \left<y,z\right> = \underbrace{\sum_{i=1}^n \left|x_i\right| \cdot 1}_{=\lVert x\rVert_1} \leq \lVert x\rVert_2 \sqrt{n} \implies \lVert x \rVert_1 \leq \lVert x_2\rVert \sqrt{n}
\]
כעת, על מנת למצוא וקטור שמוכיח שהחסם הוא הדוק, נצטרך למצוא וקטור שאצלו אי-שוויון קושי שוורץ מקיים שוויון\\
כידוע מההרצאה, נצטרך ששני הוקטורים יהיו פרופורציונליים, על מנת שוקטור יהיה פרופורציונלי ל-$z$ הוא יצטרך להיות מהצורה $\begin{pmatrix}k &\ldots& k \end{pmatrix}^t$ עבור $k\in \mathbb{R}$ כלשהו, לשם הנוחות נבחר ב-$z$ עצמו (שמקיים את התנאי עם $k=1$), ואכן
\begin{gather*}
  \lVert z \rVert_1 = \sum_{i=1}^n \left|1\right| = n\\
  \lVert z\rVert_2 = \sqrt{\sum_{i=1}^n 1^2} = \sqrt{n}
\end{gather*}
אכן קיבלנו כי
\[
  n = \sqrt{n}\sqrt{n}
\]
כנדרש
\section*{$C_{1,\infty}$}
נסמן ב-$k$ את האינדקס המקיים $x_k = \max_{i\in\left[n\right]} \left|x_i\right|$, ונשים לב כי
\[
  \lVert x\rVert_1 = \sum_{i=1}^n \left|x_i\right| \leq \sum_{i=1}^n \left|x_k\right|= \left|x_k\right| \sum_{i=1}^n 1 = n\left|x_k\right| = n \lVert x\rVert_\infty
\]
כעת נמצא את השוויון, נתבונן בוקטור $z=\begin{pmatrix}1&\ldots&1\end{pmatrix}^t$
\begin{gather*}
  \lVert z\rVert_1 = n\\
  \lVert z \rVert _\infty = \max\left\{1\right\} = 1
\end{gather*}
ואכן
\[
  n = n \cdot 1
\]
\section*{$C_{2,1}$}
נתבונן ב-$\lVert x\rVert_1^2$
\begin{gather*}
    \lVert x \rVert_1^2 = \left(\sum_{i=1}^n \left|x_i\right|\right)^2 = \sum_{i\in\left[n\right]} x_i^2 + \underbrace{\sum_{i\in\left[n\right]}\sum_{j\in\left[n\right]\setminus \left\{i\right\}} \left|x_i\right| \left|x_j\right|}_{\geq 0}
    \\ \geq \sum_{i\in\left[n\right]} x_i^2 = \left(\sqrt{\sum_{i\in\left[n\right]} x_i^2}\right)^2 = \lVert x\rVert_2^2
\end{gather*}
כעת נראה כי $e_1=\begin{pmatrix}1&0&\ldots&0\end{pmatrix}$ מקיים $\lVert e_1 \rVert_2 = \lVert e_1 \rVert_1$
\[
  \lVert e_1 \rVert_2 = \sqrt{1^2 + \sum_{i=2}^n 0^2} = \sqrt{1^2} = 1
\]
\[
  \lVert e_1 \rVert_1 = 1 + \sum_{i=2}^n \left|0\right| = 1
\]
\section*{$C_{2,\infty}$}
נסמן ב-$k$ את האינדקס המקיים $x_k = \max_{i\in\left[n\right]} \left|x_i\right|$, ונשים לב כי
\[
  \lVert x\rVert_2^2  = \sum_{i=1}^n x_i^2 \leq \sum_{i=1}^n x_k^2 = x_k^2 \sum_{i=1}^n 1^2 = x_k^2 n = \lVert x_k\rVert_\infty^2 n \implies \lVert x\rVert_2 \leq \sqrt{n} \lVert x\rVert_\infty
\]
כעת נמצא את השוויון, נתבונן בוקטור $z=\begin{pmatrix}1&\ldots&1\end{pmatrix}^t$
\begin{gather*}
  \lVert z\rVert_2 = \sqrt{\sum_{i=1}^n 1} = \sqrt{n}\\
  \lVert z \rVert_\infty = 1
\end{gather*}
ואכן
\[
  \sqrt{n} = \sqrt{n} \cdot 1
\]
\section*{$C_{\infty,1}$}
נסמן ב-$k$ את האינדקס המקיים $x_k = \max_{i\in\left[n\right]} \left|x_i\right|$, ונשים לב כי
\[
  \lVert x\rVert_\infty = \left|x_k\right| \leq  \left|x_k\right| \underbrace{\sum_{i\in\left[n\right]\setminus \left\{k\right\}} \left|x_i\right|}_{\geq 0} = \lVert x\rVert_1
\]
כעת נראה כי $e_1=\begin{pmatrix}1&0&\ldots&0\end{pmatrix}$ מקיים $\lVert e_1 \rVert_\infty = \lVert e_1 \rVert_1$
\begin{gather*}
  \lVert e_1\rVert_\infty = 1\\
  \lVert e_1\rVert_1 = 1
\end{gather*}
\section*{$C_{\infty,2}$}
נסמן ב-$k$ את האינדקס המקיים $x_k = \max_{i\in\left[n\right]} \left|x_i\right|$, ונשים לב כי
\[
  \lVert x\rVert_2^2 = \sum_{i\in\left[n\right]} x_i^2 \geq  x_k^2 = \lVert x\rVert_\infty \implies \lVert x\rVert_2 \geq \lVert x \rVert_\infty
\]

כעת נראה כי $e_1=\begin{pmatrix}1&0&\ldots&0\end{pmatrix}$ מקיים $\lVert e_1 \rVert_2 = \lVert e_1 \rVert_\infty$
\begin{gather*}
  \lVert e_1 \rVert_2 = \sqrt{1^2 \sum_{i=2}^{n} 0^2} \\
  \lVert e_1 \rVert_\infty = 1 \\
\end{gather*}
\newpage
\part*{שאלה 2}
תהיינה $\lVert \cdot \rVert_\alpha, \lVert \cdot \rVert_\beta$ שתי נורמות שקולות ב-$\mathbb{R}^n$.
\section*{סעיף א'}
יש להוכיח שכל כדור פתוח בנורמה $\lVert\cdot\rVert_\alpha$ \underline{מכיל} כדור פתוח בנורמה $\lVert\cdot\rVert_\beta$, וגם \underline{מוכל} בכדור פתוח בנורמה $\lVert\cdot\rVert_\beta$
\section*{סעיף ב'}
יש להראות שקבוצה $\Omega \subseteq \mathbb{R}^n$ היא פתוחה ביחס לנורמה $\lVert\cdot\rVert_\alpha$ אם ורק אם היא פתוחה ביחס לנורמה $\lVert\cdot\rVert_\beta$
\part*{פתרון 2}
יהיו $R\in\mathbb{R}_{>0}$ ו-$p\in\mathbb{R}^n$, נסמן את הכדור הפתוח לפי המטריקה $\lVert \cdot \rVert_\alpha$ שרדיוסו $R$ ומרכזו בנקודה $p$ בדרך הבאה
\[
  \mathcal{B}_\alpha\left(R,p\right) = \left\{x\in\mathbb{R}^n:\lVert x - p\rVert_\alpha < R\right\}
\]
משקילות הנורמות, קיים $C_{\beta,\alpha}$ כך שלכל $x\in\mathbb{R}^n$
\[
  \lVert x \rVert_\beta \leq C_{\beta,\alpha} \lVert x \rVert_\alpha
\]
נתבונן ב-$\mathcal{B}_\beta\left(\frac{1}{C_{\beta,\alpha}}\cdot R,p\right)$
\[
  \mathcal{B}_\beta\left(\frac{1}{C_{\beta,\alpha}}\cdot R,p\right)= \left\{x\in\mathbb{R}^n : \lVert x-p\rVert_\beta < \frac{1}{C_{\beta,\alpha}} \cdot R\right\}
\]
ונראה כי $\mathcal{B}_\alpha\left(R,p\right)\subseteq\mathcal{B}_\beta\left(\frac{1}{C_{\beta,\alpha}}\cdot R,p\right)$, \\
יהי $x\in \mathcal{B}_\alpha\left(R,p\right)$

\begin{gather*}
  \lVert x-p\rVert_\beta \underset{\text{שקילות הנורמות}}{\leq} C_{\beta,\alpha} \lVert x-p \rVert_\alpha \underset{x\in\mathcal{B}_\alpha\left(R,p\right)}{<} C_{\beta,\alpha}R\\
  \implies \lVert x-p \rVert_\beta<\frac{1}{C_{\beta,\alpha}} R
\end{gather*}
באותה מידה, קיים גם $C_{\alpha,\beta}$, כך שלכל $x\in\mathbb{R}^n$
\[
  \lVert x\rVert_\alpha \leq C_{\alpha,\beta} \lVert x\rVert_\beta
\]
כעת נראה כי 
\[
  \mathcal{B}_\alpha \left(R,p\right) \supseteq \mathcal{B}_\beta \left(\frac{1}{C_{\alpha,\beta}}\cdot R, p\right) = \left\{x\in\mathbb{R}^n :\lVert x-p\rVert_\beta < \frac{1}{C_{\alpha,\beta}} \cdot R\right\}
\]
יהי $x\in\mathcal{B}_\beta\left(\frac{1}{C_{\alpha,\beta}},p\right)$
\begin{gather*}
  \lVert x-p\rVert_\alpha \underset{\text{שקילות הנורמות}}{\leq} C_{\alpha,\beta} \lVert x-p \rVert_\beta \underset{x\in\mathcal{B}_\beta\left(\frac{1}{C_{\alpha,\beta}},p\right)}{<} R\\
  \implies \lVert x-p \rVert_\alpha< R
\end{gather*}
\section*{סעיף ב'}
תהי $\Omega \subseteq \mathbb{R}^n$ ו-$\lVert \cdot \rVert_\alpha, \lVert \cdot \rVert_\beta$ שתי נורמות
\subsection*{$\impliedby$}
נניח כי $\Omega$ פתוחה ביחס ל-$\lVert \cdot \rVert_\alpha$\\
יהי $x\in\Omega$, מההנחה קיים $\varepsilon>0$ כך ש-$x\in\mathcal{B}_\alpha\left(\varepsilon,x\right)\subseteq \Omega$, מסעיף א', קיים $C_{\alpha,\beta}$ כך ש-$x\in\mathcal{B}_\beta\left(\frac{1}{C_{\alpha,\beta}}\cdot\varepsilon,x\right)\subseteq \mathcal{B}_\alpha\left(\varepsilon,x\right)\subseteq \Omega$\\
$\mathcal{B}_\beta\left(\frac{1}{C_{\alpha,\beta}}\cdot\varepsilon,x\right)$ היא סביבה פתוחה לפי $\lVert \cdot \rVert_\beta$, אשר מכילה את $x$, ומוכלת בתוך $\Omega$\\
משרירותיות $x$, קיבלנו כי לכל $x\in\Omega$, קיימת סביבה פתוחה לפי $\lVert \cdot \rVert_\beta$ המוכלת כולה ב-$\Omega$ ומכילה את $x$, לכן לפי ההגדרה $\Omega$ פתוחה לפי $\lVert \cdot \rVert_\beta$
\subsection*{$\implies$}
כעת נניח כי $\Omega$ פתוחה ביחס ל-$\lVert\cdot\rVert_\beta$\\
יהי $x\in\Omega$, מההנחה קיים $\varepsilon>0$ כך ש-$x\in\mathcal{B}_\beta\left(\varepsilon,x\right)\subseteq \Omega$, מסעיף א' קיים גם $C_{\beta,\alpha}$ כך ש $x\in\mathcal{B}_\alpha\left(\frac{1}{C_{\beta,\alpha}}\cdot\varepsilon,x\right)\subseteq\mathcal{B}_\beta\left(\varepsilon,x\right)\subseteq\Omega$\\
$\mathcal{B}_\alpha\left(\frac{1}{C_{\beta,\alpha}}\cdot\varepsilon,x\right)$ היא סביבה פתוחה לפי $\lVert \cdot \rVert_\beta$ אשר מכילה את $x$, ומוכלת בתוך $\Omega$\\
משרירותיות $x$, קיבלנו כי לכל $x\in\Omega$ קיימת סביבה פתוחה לפי $\lVert \cdot \rVert_\alpha$ המוכלת כולה ב-$\Omega$ ומכילה את $x$, לכן לפי ההגדרה $\Omega$ פתוחה לפי $\lVert \cdot \rVert_\alpha$
\newpage
\part*{שאלה 3}
תהי $A:\mathbb{R}^n\to\mathbb{R}^m$ טרנספורמציה לינארית, נסמן 
\[
  \lVert A\rVert = \sup_{\lVert v\rVert=1} \left\{\lVert Av\rVert\right\}
\]
\section*{סעיף א'}
יש להוכיח כי $\lVert A\rVert$ נורמה
\section*{סעיף ב'}
עבור נורמות הערך המוחלט יש להוכיח כי
\[
  \lVert A\rVert = \max_{j\in\left[n\right]}\left\{ \sum_{i=1}^m \left|a_{i,j}\right|\right\}
\]
\part*{פתרון 3}
\section*{סעיף א'}
נראה כי $\lVert A \rVert$ מקיימת את שלושת התכונות של נורמה
\subsection*{חיוביות}
ראשית נראה כי $\lVert A\rVert \geq 0$\\
$\lVert A\rVert$ מוגדר להיות הסופרימום של קבוצת נורמות כלשהן, מתכונת הנורמה $\lVert Av\rVert$ ב-$\mathbb{R}^m$, $\lVert Av\rVert\geq 0$ לכל $v$ (ובפרט לכל $v$ במקיים $\rVert v\rVert=1$)\\
כעת נראה ש-$\lVert A\rVert=0 \iff A=0_{m\times n}$\\
על מנת שהסופרימום של קבוצת אי-שליליים יהיה $0$, כל איברי הקבוצה צריכים להיות $0$.\\
כלומר
\[
  \left(\forall v\in\mathbb{R}^n\right)\left(\lVert v\rVert=1\implies \lVert Av\rVert = 0\right)
\]
מכיוון ש-$\lVert Av\rVert$ היא נורמה, $\lVert Av\rVert = 0 \iff Av = 0_{m}$\\
הטרנספורמציה \underline{היחידה} המקיימת את זה היא טרנספורמציית האפס
\subsection*{הומוגניות}
יהי $\alpha \in \mathbb{R}$, נרצה להראות כי $\lVert \alpha A\rVert=\left|\alpha\right|\lVert A\rVert$
\[
  \lVert \alpha A\rVert = \sup_{\lVert v\rVert=1}\left\{\lVert \alpha A v\rVert\right\} \underset{\left(*\right)}{=} \sup_{\lVert v\rVert =1} \left\{\left|\alpha\right|\lVert Av\rVert\right\} \underset{\left(**\right)}{=} \left|\alpha\right|\sup_{\lVert v\rVert =1} \left\{\lVert Av\rVert\right\} 
\]
$\left(*\right)$ $\lVert Av\rVert$ היא נורמה ולכן מקיימת הומוגניות\\
$\left(**\right)$ כפל של כל איבר בקבוצה בסקלאר חיובי
\subsection*{אי שוויון המשולש}
נרצה להוכיח כי
\[
  \left(\forall B\in M_{m \times n}\left(\mathbb{R}\right)\right)\left(\lVert A + B\rVert\leq\lVert A\rVert + \lVert B \rVert\right)
\]
\begin{gather*}
  \lVert A+B\rVert = \sup_{\lVert v\rVert = 1} \left\{\lVert \left(A+B\right) v\rVert\right\} = \sup_{\lVert v\rVert = 1} \left\{\lVert \left(Av\right) +\left(Bv\right)\rVert\right\} \\
  \underset{\left(*\right)}{\leq} \sup_{\lVert v\rVert = 1} \left\{\lVert A v\rVert +\lVert B v\rVert\right\}  \underset{\left(**\right)}{\leq} \sup_{\lVert v\rVert = 1}\left\{\lVert Av\rVert\right\} + \sup_{\lVert v\rVert = 1}\left\{\lVert Bv\rVert\right\} = \lVert A\rVert + \lVert B\rVert
\end{gather*}
$\left(*\right)$ אי-שוויון המשולש של $\lVert Av\rVert$\\
$\left(**\right)$ "בחירה" נפרדת של $v$ עבור $A$ ועבור $B$ יכולה רק להגדיל את הסופרימום
קיבלנו כי $\lVert A \rVert$ מקיימת את שלושת התכונות של נורמה, ולכן היא נורמה
\section*{סעיף ב'}
נוכיח על ידי שני אי-שוויונות, נגדיר $M=\max_{j\in\left[n\right]}\left\{\sum_{i=1}^m\left|a_{i,j}\right|\right\}$.\\
נוכיח כי $\lVert A\rVert \leq M$, ואז נוכיח כי $\lVert A\rVert \geq M$
\subsection*{$\lVert A\rVert \leq M$}
עבור $v=\begin{pmatrix}v_1\\\vdots\\v_n\end{pmatrix}$ המקיים $\sum_{j=1}^{n} \left|v_j\right|=1$,
נכתוב ישירות את $\lVert Av\rVert$
\begin{gather*}
  \lVert Av \rVert = \sum_{i=1}^m\left| \sum_{j=1}^n a_{i,j} v_j\right| \underset{\triangle}{\leq}\sum_{i=1}^m\sum_{j=1}^n \left|a_{i,j} v_j\right|\\
  =\sum_{i=1}^m\sum_{j=1}^n \left|a_{i,j}\right|\left| v_j\right| \underset{\left(*\right)}{=} \sum_{j=1}^n\left(\sum_{i=1}^m \left|a_{i,j}\right|\right)\left| v_j\right| \\
  \underset{\left(**\right)}{\leq} \sum_{j=1}^n M \left|v_j\right| = M \sum_{j=1}^n \left|v_j\right| \underset{\lVert v \rVert = 1}{=} M
\end{gather*}
$\left(*\right)$ קומוטטיביות בסכום סופי\\
$\left(**\right)$  לכל $j\in\left[n\right]$, $\sum_{i=1}^m \left|a_{i,j}\right|\leq \max_{k\in\left[n\right]} \sum_{i=1}^m \left|a_{i,k}\right|=M$\\
מכיוון שלכל $v$ כנ"ל מתקיים $\lVert Av \rVert \leq M$, מתקיים שגם הסופרימום חסום ע"י $M$, כלומר
\[
  \lVert A \rVert = \sup_{\lVert v \rVert = 1}\left\{\lVert Av\rVert\right\}
\]
\subsection*{$\lVert A \rVert \geq M$}
נגדיר את $k\in\left[n\right]$ להיות האינדקס שבו $\sum^m_{i=1}\left|a_{i,k}\right|$ מקבל ערך מקסימלי.\\
מהתנאי ש-$\lVert v\rVert=1$, צריך שיתקיים $\sum_{j=1}^n\left|v_j\right|=1$.\\
נרצה להגדיל את $\lVert Av\rVert$, נרצה ש-$\left|v_k\right|$ יהיה גדול ככל האפשר, זה מתקבל עבור $v=e_k$
\end{document}
